% GENERATED FILE. Edit canonical lessons, then run npm run generate:books.
% canonical-source-hash: fnv1a64:0a2a2d1b495caa10
% canonical-lessons: KA-C91-address, KA-C91-birthday, KA-C91-wedding, KA-C91-hometown, KA-C91-language

\chapter{About You}
\label{ch:ka-about-you}
\hlchaptermodality{91}

\begin{chapteropening}
\textbf{By the end of this chapter:} \emph{I can talk about my address, a birthday, a wedding, my home town and a language.}

Complete the last lesson of chapter 91: i can talk about my address, a birthday, a wedding, my home town and a language.
\end{chapteropening}

\section[viḷāsa]{\kn{ವಿಳಾಸ} (viḷāsa) — address}

\label{lesson:KA-C91-address}



\begin{quote}
Before the new one: say the Kannada for the neighbours, then the Kannada for uncle.
\end{quote}

\subsection*{You'll want to know: \kn{ವಿಳಾಸ}}
\textbf{\kn{ವಿಳಾಸ}} (\emph{viḷāsa}) — \textquotedblleft{}address\textquotedblright{}.

\textbf{\kn{ನಿಮ್ಮ} \kn{ವಿಳಾಸ} \kn{ಏನು}?} (\emph{nimma viḷāsa ēnu?}) — \textquotedblleft{}what is your address?\textquotedblright{}, built like \textbf{\kn{ನಿಮ್ಮ} \kn{ಹೆಸರು} \kn{ಏನು}?} from the start of the book.

\begin{grammarlens}[title={What you've built}]
The name question, asked about a place.
\end{grammarlens}

\subsection*{Your turn}
\begin{itemize}
\raggedright
  \item \emph{Say it:} \emph{viḷāsa}
  \item \emph{Say it:} \emph{viḷāsa}, once more
  \item \emph{Say it:} ask \emph{nimma viḷāsa ēnu?}
  \item \emph{Recall:} say the Kannada for the neighbours, then the Kannada for uncle, then say \emph{viḷāsa} again
\end{itemize}

\begin{tcolorbox}[breakable,colback=teal!4,colframe=teal!35!black,title={Before you move on}]
What does \kn{ವಿಳಾಸ} mean? (\textquotedblleft{}Address\textquotedblright{}.) Say it once more.
\end{tcolorbox}
\section[huṭṭuhabba]{\kn{ಹುಟ್ಟುಹಬ್ಬ} (huṭṭuhabba) — birthday}

\label{lesson:KA-C91-birthday}



\begin{quote}
Before the new one: say the Kannada for uncle, then the Kannada for address.
\end{quote}

\subsection*{You'll want to know: \kn{ಹುಟ್ಟುಹಬ್ಬ}}
\textbf{\kn{ಹುಟ್ಟುಹಬ್ಬ}} (\emph{huṭṭuhabba}) — \textquotedblleft{}birthday\textquotedblright{}, literally \textquotedblleft{}birth festival\textquotedblright{}: \textbf{\kn{ಹುಟ್ಟು}} (\emph{huṭṭu}), \textquotedblleft{}birth\textquotedblright{}, and \emph{habba}, \textquotedblleft{}festival\textquotedblright{}, which gets a lesson of its own later.

\begin{grammarlens}[title={What you've built}]
A festival of one's own.
\end{grammarlens}

\subsection*{Your turn}
\begin{itemize}
\raggedright
  \item \emph{Say it:} \emph{huṭṭuhabba}
  \item \emph{Say it:} \emph{huṭṭuhabba}, once more
  \item \emph{Say it:} say \emph{huṭṭuhabba}
  \item \emph{Recall:} say the Kannada for uncle, then the Kannada for address, then say \emph{huṭṭuhabba} again
\end{itemize}

\begin{tcolorbox}[breakable,colback=teal!4,colframe=teal!35!black,title={Before you move on}]
What does \kn{ಹುಟ್ಟುಹಬ್ಬ} mean? (\textquotedblleft{}Birthday\textquotedblright{}.) Say it once more.
\end{tcolorbox}
\section[maduve]{\kn{ಮದುವೆ} (maduve) — wedding, marriage}

\label{lesson:KA-C91-wedding}



\begin{quote}
Before the new one: say the Kannada for address, then the Kannada for birthday.
\end{quote}

\subsection*{You'll want to know: \kn{ಮದುವೆ}}
\textbf{\kn{ಮದುವೆ}} (\emph{maduve}) — \textquotedblleft{}wedding, marriage\textquotedblright{}.

\textbf{\kn{ಮದುವೆಗೆ} \kn{ಆಹ್ವಾನ}} (\emph{maduvege āhvāna}) — \textquotedblleft{}an invitation to a wedding\textquotedblright{}, with \textbf{\kn{ಆಹ್ವಾನ}}, \textquotedblleft{}invitation\textquotedblright{}, from the last tranche.

\begin{grammarlens}[title={What you've built}]
The biggest invitation of all.
\end{grammarlens}

\subsection*{Your turn}
\begin{itemize}
\raggedright
  \item \emph{Say it:} \emph{maduve}
  \item \emph{Say it:} \emph{maduve}, once more
  \item \emph{Say it:} say \emph{maduvege āhvāna}
  \item \emph{Recall:} say the Kannada for address, then the Kannada for birthday, then say \emph{maduve} again
\end{itemize}

\begin{tcolorbox}[breakable,colback=teal!4,colframe=teal!35!black,title={Before you move on}]
What does \kn{ಮದುವೆ} mean? (\textquotedblleft{}Wedding, marriage\textquotedblright{}.) Say it once more.
\end{tcolorbox}
\section[ūru]{\kn{ಊರು} (ūru) — town, home town}

\label{lesson:KA-C91-hometown}



\begin{quote}
Before the new one: say the Kannada for birthday, then the Kannada for a wedding.
\end{quote}

\subsection*{You'll want to know: \kn{ಊರು}}
\textbf{\kn{ಊರು}} (\emph{ūru}) — \textquotedblleft{}town\textquotedblright{}, and in particular one's \textbf{home town}. It is the \emph{-ur} at the end of place names across the south: Bengaluru, Mysuru.

\textbf{\kn{ನಿಮ್ಮ} \kn{ಊರು} \kn{ಯಾವುದು}?} (\emph{nimma ūru yāvudu?}) — \textquotedblleft{}where is your home town?\textquotedblright{}

\begin{grammarlens}[title={What you've built}]
The word hiding in every -uru city name.
\end{grammarlens}

\subsection*{Your turn}
\begin{itemize}
\raggedright
  \item \emph{Say it:} \emph{ūru}
  \item \emph{Say it:} \emph{ūru}, once more
  \item \emph{Say it:} ask \emph{nimma ūru yāvudu?}
  \item \emph{Recall:} say the Kannada for birthday, then the Kannada for a wedding, then say \emph{ūru} again
\end{itemize}

\begin{tcolorbox}[breakable,colback=teal!4,colframe=teal!35!black,title={Before you move on}]
What does \kn{ಊರು} mean? (\textquotedblleft{}Town, home town\textquotedblright{}.) Say it once more.
\end{tcolorbox}
\section[bhāṣe]{\kn{ಭಾಷೆ} (bhāṣe) — language}

\label{lesson:KA-C91-language}



\begin{quote}
Before the new one: say the Kannada for a wedding, then the Kannada for home town.
\end{quote}

\subsection*{You'll want to know: \kn{ಭಾಷೆ}}
\textbf{\kn{ಭಾಷೆ}} (\emph{bhāṣe}) — \textquotedblleft{}language\textquotedblright{}.

\textbf{\kn{ಕನ್ನಡ} \kn{ಭಾಷೆ}} (\emph{kannaḍa bhāṣe}) — \textquotedblleft{}the Kannada language\textquotedblright{}.

\begin{grammarlens}[title={What you've built}]
The word for what you are learning.
\end{grammarlens}

\subsection*{Your turn}
\begin{itemize}
\raggedright
  \item \emph{Say it:} \emph{bhāṣe}
  \item \emph{Say it:} \emph{bhāṣe}, once more
  \item \emph{Say it:} say \emph{kannaḍa bhāṣe}
  \item \emph{Recall:} say the Kannada for a wedding, then the Kannada for home town, then say \emph{bhāṣe} again
\end{itemize}

\begin{tcolorbox}[breakable,colback=teal!4,colframe=teal!35!black,title={Before you move on}]
What does \kn{ಭಾಷೆ} mean? (\textquotedblleft{}Language\textquotedblright{}.) Say it once more.
\end{tcolorbox}
